\begin{flushright}
  \fechaclase{15 de octubre de 2026}
\end{flushright}

\section{Control deslizante basada en el método de la elipsoide atractiva (SMC-AEM)}

Sea el sistema con perturbaciones mismatched

\begin{equation*}
    \begin{aligned}
        \dot{x}_1 &= f_m(x,t) + x_2 \\
        \dot{x}_2 &= f_u(x, u, t) + u \\
        y &= x
    \end{aligned}
\end{equation*}

con \(|f_1|\leq D_m\) y \(|f_2|\leq D_u\). Se desea realizar seguimineto de trayectorias para la referencia r. Se tiene el error \(e_1=r-x_1\) y \(e_2=\dot{r}-x_2\) con derivadas:

\begin{equation*}
    \begin{aligned}
        \dot{e}_1 &= \dot{r}-\dot{x}_1 = \dot{r} - f_m(x,t) - x_2 = e_2 - f_m(x,t)\\
        \dot{e}_2 &= \ddot{r}-\dot{x}_2 = \ddot{r} - f_u(x, t) - u
    \end{aligned}
\end{equation*}

Se propone la superficie

\begin{equation*}
    s = e_2 + me_1 \quad, m>0
\end{equation*}

Calculamos la derivada de la superficie

\begin{equation*}
    \begin{aligned}
        \dot{s} &= \dot{e}_2 + m\dot{e}_1 = \ddot{r} - f_u(x, t) - u + m(e_2 - f_m(x,t))\\
        &= \ddot{r} + me_2 - u - mf_u(x,t) - f_m(x, t)
    \end{aligned}
\end{equation*}

proponemos la función de Lyapunov

\begin{equation*}
    V = \frac{1}{2}s^2
\end{equation*}

derivando

\begin{equation*}
    \dot{V} = s\dot{s} = s(\ddot{r} + m(e_2+D_u)+D_m -u)
\end{equation*}

con el controlador

\begin{equation*}
    U = me_2 + k\text{sign}(s) 
\end{equation*}

\begin{equation*}
    V_s \leq s (\ddot{r} + mD_u + D_m -k\text{sign}(s)) \leq (R_2+mD_u+D_m)|s|-K|s|
\end{equation*}

elegimis \(K>R_2+mD_u+D_m\)

\begin{equation*}
    \dot{V} \leq -(K-R_2-mD_u-D_m)
\end{equation*}

esto garantiza que

\begin{equation*}
    t_r \leq \frac{|s(0)|}{\eta}
\end{equation*}

con \(\eta=K-R_2-mD_u-D_m\). Cuando $s=0$ , $e_2+me_1=0$, esto es

\begin{equation*}
    \dot{e}_1 = - me_1 + D_u
\end{equation*}

Se define la función de Lyapunov

\begin{equation*}
    V_1 = pe_1^2
\end{equation*}

con derivada

\begin{equation*}
    \begin{aligned}
        \dot{V}_1 &= 2pe_1\dot{e}_1\\
        &=2pe_1(-me_1+D_u)\\
        &=-2pme_1^2 + 2pe_1f_u
    \end{aligned}
\end{equation*}

usando la desigundad de Young

\begin{equation*}
    2pe_1du \leq p\epsilon e_1^2 + \frac{p}{\epsilon}f_u^2 \leq p\epsilon e_1^2 + \frac{p}{\epsilon}D_u^2
\end{equation*}

entonces

\begin{equation*}
    \begin{aligned}
        \dot{V}_1 \leq -2pme_1^2 + p\epsilon e_1^2 + \frac{p}{\epsilon}D_u^2 \\
        \leq p(-2m+\epsilon)e_1^2 + \frac{p}{\epsilon}D_u^2
    \end{aligned}
\end{equation*}

deseamos que 

\begin{equation*}
    \dot{V} \leq -\alpha V + \beta
\end{equation*}

demostrnaod que

\begin{equation*}
    -2m+\epsilon \leq - \alpha
\end{equation*}

entonces 

\begin{equation*}
    \dot{V} \leq - \alpha pe_1^2+\beta= - \alpha V+ \beta
\end{equation*}

lo que significa que 

\begin{equation*}
    V \leq \frac{\beta}{\alpha} \quad \text{o bien} \quad |e_1| \leq \sqrt{\frac{\beta}{\alpha p}}
\end{equation*}

\begin{equation*}
    |e_1| \leq \frac{D_u}{\sqrt{\alpha\epsilon}}
\end{equation*}

asi el radio de la elipsoide es 

\begin{equation*}
    r= \frac{D_u}{\sqrt{\alpha\epsilon}}
\end{equation*}

con 

\begin{equation*}
    \alpha = 2m - \epsilon
\end{equation*}